% TOP_PROBLEM: {"id": 30007284, "title": "Open homomorphisms of solvably closed global Galois groups", "category_id": 1, "category": {"id": 1, "name": "number_theory", "display_name": "Number Theory", "description": "Properties of integers, prime numbers, Diophantine equations.", "slug": "number-theory", "order_index": 1, "created_at": "2026-07-31T15:26:25.670Z"}, "status": "open"}
% FIRST_POSTED: 2026-10-06
% ATTEMPT_STATUS: solved
% OPENAI_CLAIM: {"schema_version":1,"source_repo":"https://github.com/openai/math","source_commit":"adc7f1241b42e322a6451854ab7e4b4c146bf78a","release_date":"2026-10-06","match":"full","resolution":"proved","independently_reviewed":false,"second_pass":"agreed","families":[{"family":"031","title":"Uchida’s conjecture for open homomorphisms of Galois groups","subject":"Number theory","manuscripts":[{"title":"Open Homomorphisms of Global Solvably Closed Galois Groups","pdf_path":"preprints/Open-Homomorphisms-of-Global-Solvably-Closed-Galois-Groups-October-5-2026/open-homomorphisms-solvably-closed-galois-groups.pdf","readme_path":"preprints/Open-Homomorphisms-of-Global-Solvably-Closed-Galois-Groups-October-5-2026/README.md","date":"2026-10-05","theorem_ref":"thm:main (sections/introduction.tex:11)","statement_sources":[{"path":"preprints/Open-Homomorphisms-of-Global-Solvably-Closed-Galois-Groups-October-5-2026/build/sections/introduction.tex","line":11,"label":"thm:main"}]}],"lean":null,"reasoning_trace":null}]}
\documentclass[11pt]{article}
\usepackage[margin=25mm]{geometry}
\usepackage{fontspec}
\setmainfont{Latin Modern Roman}
\usepackage{amsmath,amssymb,mathtools}
\usepackage{unicode-math}
\setmathfont{Latin Modern Math}
\usepackage{enumitem}
\usepackage{xurl,hyperref,bookmark}
\hypersetup{colorlinks=true,urlcolor=blue}
\setcounter{secnumdepth}{0}
\setlength{\parindent}{0pt}
\setlength{\parskip}{5pt}
\setlength{\emergencystretch}{3em}
\begin{document}
\section*{Open homomorphisms of solvably closed global Galois groups}

\subsection{OpenAI's claim}
OpenAI presents the following claims relevant to this question, with manuscript numbers following paper order in Sources. In manuscript 1 (family 031; 2026-10-05), the asserted result is that every continuous open homomorphism between the stated solvably closed number-field Galois groups arises contravariantly from a unique equivariant field embedding, with no additional kernel restriction.

\subsection{Mathematical statement of the claimed result}
\paragraph{Manuscript 1.}
Let $F_1,F_2$ be number fields and $E_i/F_i$ possibly infinite Galois extensions, with each $E_i$ solvably closed, meaning it has no nontrivial abelian algebraic extension. Every continuous open homomorphism $\alpha:\operatorname{Gal}(E_1/F_1)\to\operatorname{Gal}(E_2/F_2)$ is claimed induced by exactly one field embedding $j:E_2\hookrightarrow E_1$, characterized by
\[
 g\circ j=j\circ\alpha(g)\quad(g\in\operatorname{Gal}(E_1/F_1)).
\]
In particular $j(F_2)\subseteq F_1$; finite extensions and surjectivity of $\alpha$ are not required.

\subsection{Outline of the argument}
For manuscript 1 (family 031; 2026-10-05), power-residue separation matches local characteristics and bounds missed primes. Regular-module Kummer extensions supply enough characters to detect Frobenius images, while averaging and corestriction turn the partial local correspondence into the required field map.

\subsection{Scope relative to this problem}
The result supplies the conclusion in this problem with its displayed hypotheses and quantifiers. The broader statements and applications in the manuscript are kept separate from the particular question above. The supporting result is manuscript 1 (family 031; 2026-10-05).

\subsection{Formal verification}
No Lean formalization is listed at this snapshot for manuscript 1 (family 031; 2026-10-05).

This site has not executed the Lean code. Described statement coverage is not an independent proof check.

\subsection{Status caveat}
The manuscripts were produced by an internal OpenAI model. OpenAI cautions that results without formalization can contain mistakes. This site has not independently checked the proofs; the label records the scope of OpenAI's claim.

\subsection{Sources}
\begin{itemize}
\item OpenAI, \emph{Open Homomorphisms of Global Solvably Closed Galois Groups}, 2026-10-05; citation key \texttt{OAI:Open-Homomorphisms-of-Global-Solvably-Closed-Galois-Groups-October-5-2026}. \href{https://github.com/openai/math/blob/adc7f1241b42e322a6451854ab7e4b4c146bf78a/preprints/Open-Homomorphisms-of-Global-Solvably-Closed-Galois-Groups-October-5-2026/open-homomorphisms-solvably-closed-galois-groups.pdf}{PDF (GitHub)}; \href{https://github.com/openai/math/raw/adc7f1241b42e322a6451854ab7e4b4c146bf78a/preprints/Open-Homomorphisms-of-Global-Solvably-Closed-Galois-Groups-October-5-2026/open-homomorphisms-solvably-closed-galois-groups.pdf}{PDF (direct)}; \href{https://github.com/openai/math/blob/adc7f1241b42e322a6451854ab7e4b4c146bf78a/preprints/Open-Homomorphisms-of-Global-Solvably-Closed-Galois-Groups-October-5-2026/README.md}{Manuscript page and citation}; \href{https://github.com/openai/math/blob/main/preprints/Open-Homomorphisms-of-Global-Solvably-Closed-Galois-Groups-October-5-2026/open-homomorphisms-solvably-closed-galois-groups.pdf}{Latest version}; Statement source: \href{https://github.com/openai/math/blob/adc7f1241b42e322a6451854ab7e4b4c146bf78a/preprints/Open-Homomorphisms-of-Global-Solvably-Closed-Galois-Groups-October-5-2026/build/sections/introduction.tex\#L11}{thm:main}.
\item \href{https://github.com/openai/math/blob/adc7f1241b42e322a6451854ab7e4b4c146bf78a/CONTENTS.md}{OpenAI catalogue entry}.
\end{itemize}
\end{document}
